% Reconstructed analytical author extension after the runtime reset.
% This source has a new identity. It is not a byte recovery of the lost
% reviewed source 43e679e564e2dfc97083a779ca4c02bfad92bab8003068ada87d8d55f6054d92.
% No frozen algorithm, sample, actor, inference rule, or outcome is changed.
\section{From action safeguards to gains in the same recursion}
\label{supp:dynamicgain122}

The paid-support safeguard compares two actions while retaining a specified
future policy. Extending the economic gain to several changed dates requires
an additional recursion. The physical transition law alone cannot propagate
that gain under a nonlinear risk preference. We give the appropriate lower
recursion, retaining the original state, feasible policy, conditional risk
primitive and information set. This is an action-gain application of the
reference-centered representation in Section~\ref{sec:centeredrisk92} and of
risk-averse dynamic programming \citep{ruszczynski2010risk,ruszczynski2014erratum}.
It is not a new general policy-iteration principle, or a numerical result for
a longer horizon than the executed investment experiment.

\subsection{Original-policy advantages and centered lower evaluations}

At dates $t=0,\ldots,T$, let $\mathcal X_t$ be measurable state spaces.
For each complete feasible measurable policy $\zeta$, define its original
cost by
\begin{equation}\label{eq:policyrecursion122}
 V_T^\zeta=g,\qquad
 V_t^\zeta(x)=c_t(x,\zeta_t(x))+
       \beta_t\rho_{t,x,\zeta_t(x)}(V_{t+1}^\zeta),
 \qquad \beta_t\ge0.
\end{equation}
The terminal cost and the selected stage costs are bounded and measurable.
For every selected action, the original conditional evaluation is finite on
bounded measurable functions, is monotone, cash invariant and normalized at
zero, and maps bounded measurable continuations to bounded measurable
state functions. These hypotheses ensure that the finite recursions below
are defined. They also specify the economic preference; replacing a static
risk objective by a nested one is not allowed by these hypotheses.

Suppose the selected original risk primitives have the dual representation
\eqref{eq:dualrisk92} on this same function domain. The admissible $q$ are
probability measures, their penalties $\psi(q)$ may take the value $+\infty$,
and the set of finite-penalty measures is nonempty. All integrals below are
with respect to these original successor-state measures. Attainment of a
risk supremum is not assumed.

Fix an incumbent $\pi$ and a proposed complete feasible policy $\pi'$.
The original incumbent continuation, not a newly fitted or altered
objective, is used in the one-step advantage
\begin{equation}\label{eq:incumbentadvantage122}
 b_t^*(x)=V_t^\pi(x)-c_t(x,\pi'_t(x))
       -\beta_t\rho_{t,x,\pi'_t(x)}(V_{t+1}^\pi).
\end{equation}
Assume a bounded measurable certificate supplies
\begin{equation}\label{eq:advantagelower122}
                 0\le b_t(x)\le b_t^*(x)
                 \quad\hbox{at every state and date.}
\end{equation}
It suffices to impose this on an explicitly specified invariant reachable
state domain supporting every relevant original risk kernel. A finite
sample check alone is not this all-state hypothesis. Retaining the incumbent
at a state gives $b_t^*(x)=b_t(x)=0$ exactly.

For $v=V_{t+1}^\pi$, center at the \emph{updated} action $a=\pi'_t(x)$:
\begin{align}
 k_t^\pi(x,q)
  &=\rho_{t,x,a}(v)-\int v\,dq+\psi_{t,x,a}(q),\label{eq:gainpenalty122}\\
 \mathcal L_t^\pi h(x)
  &=\rho_{t,x,a}(v)-\rho_{t,x,a}(v-h)
    =\inf_{q\in\mathcal Q_{t,x,a}}
           \left\{\int h\,dq+k_t^\pi(x,q)\right\}.
       \label{eq:lowercenter122}
\end{align}
The equality is algebraic: subtract the risk supremum at $v-h$ from the
finite risk value at $v$. In particular, this is an infimum with a
\emph{positive} penalty. Using the supremum operator for an upper error
account here would reverse the economic guarantee.

\begin{lemma}[The lower centered operator]
\label{lem:lowercenter122}
The penalties in \eqref{eq:gainpenalty122} are nonnegative and have
infimum zero. The operator $\mathcal L_t^\pi$ is monotone, cash invariant,
concave, normalized at zero and nonexpansive in the supremum norm.
If verified penalties satisfy
\[
        0\le\underline k_t(x,q)\le k_t^\pi(x,q)
                    \quad(q\in\mathcal Q_{t,x,\pi'_t(x)}),
\]
then the operator
\[
 \underline{\mathcal L}_t h(x)
   =\inf_{q\in\mathcal Q_{t,x,\pi'_t(x)}}
            \left\{\int h\,dq+\underline k_t(x,q)\right\}
\]
has these same properties and satisfies
$\underline{\mathcal L}_t h\le\mathcal L_t^\pi h$.
For this second operator, measurability of its statewise outputs is also
assumed. Zero penalties are an allowed conservative choice.
\end{lemma}

\begin{proof}
Each feasible dual objective at $v$ is bounded above by its supremum,
so $k_t^\pi\ge0$. Approximation of that supremum by feasible measures
gives $\inf_q k_t^\pi=0$, including when no maximizing measure exists.
Since $0\le\underline k_t\le k_t^\pi$, the lower penalties also have
infimum zero. An infimum of affine functions of $h$ is concave.
Probability measures preserve order and integrate a constant to that
constant; taking infima proves monotonicity and cash invariance.
The zero infimum of the penalties proves normalization. If
$\|h-\widetilde h\|_\infty\le\epsilon$, monotonicity and cash invariance
give
$\mathcal L_t^\pi\widetilde h-\epsilon
 \le\mathcal L_t^\pi h
 \le\mathcal L_t^\pi\widetilde h+\epsilon$,
and the same proof applies to $\underline{\mathcal L}_t$.
Taking infima after the pointwise penalty inequality gives the claimed
operator order. These properties also show that each output lies between
$-\|h\|_\infty$ and $\|h\|_\infty$; with the stated measurability
hypothesis it is a bounded measurable state function.
\end{proof}

\begin{theorem}[Propagation of certified action gains]
\label{thm:dynamicgain122}
Under the preceding assumptions, let
$D_t=V_t^\pi-V_t^{\pi'}$. Define the exact and conservative lower-gain
recursions by
\begin{align}
 H_T=\underline H_T&=0,\nonumber\\
 H_t&=b_t+\beta_t\mathcal L_t^\pi H_{t+1},\label{eq:gainrecursion122}\\
 \underline H_t&=b_t+\beta_t
                     \underline{\mathcal L}_t\underline H_{t+1}.
                    \nonumber
\end{align}
Then, at every state and date,
\begin{equation}\label{eq:dynamicgainbound122}
          0\le\underline H_t\le H_t\le D_t,
 \qquad V_t^{\pi'}\le V_t^\pi.
\end{equation}
Moreover, any bounded measurable nonnegative functions $h_t$ with
$h_T=0$ and
\[
 h_t\le b_t+\beta_t\underline{\mathcal L}_t h_{t+1}
\]
are valid lower certificates: $h_t\le\underline H_t\le D_t$.
Thus a verified subsolution can certify the gain without solving the
lower recursion exactly.
\end{theorem}

\begin{proof}
Subtract the two original-policy recursions, add and subtract
$c_t(x,\pi'_t(x))+
 \beta_t\rho_{t,x,\pi'_t(x)}(V_{t+1}^\pi)$,
and use \eqref{eq:lowercenter122}. This gives the exact identity
\begin{equation}\label{eq:exactdynamicgain122}
        D_T=0,\qquad
        D_t=b_t^*+\beta_t\mathcal L_t^\pi D_{t+1}.
\end{equation}
The terminal equality uses the same original terminal cost for both
policies. If $D_{t+1}\ge H_{t+1}\ge\underline H_{t+1}\ge0$,
monotonicity, the nonnegative discount and $b_t^*\ge b_t\ge0$
give $D_t\ge H_t\ge\underline H_t\ge0$.
Backward induction from the terminal equality proves
\eqref{eq:dynamicgainbound122}. The same induction using the subsolution
inequality gives $h_t\le\underline H_t$. No physical expectation,
fixed adverse kernel, change of risk preference, or equality of
successive policies has been inserted in this calculation.
\end{proof}

The theorem supplies the missing link between a safe action update and
its economic effect through a common recursion. If only a single date
changes, its advantage is propagated through the unchanged earlier risk
operators. If several dates change, each advantage is still computed
against the same incumbent continuation sequence; evaluating datewise
advantages under unrelated fitted futures would not prove the result.
Repeated complete-policy updates can be certified by applying the theorem
anew to each incumbent. For finitely many such updates, the sum of their
statewise lower-gain certificates is a lower bound on the total gain,
because the actual original values telescope. This is not a convergence
rate to the globally optimal policy.

\subsection{A finite occupation calculation with the correct orientation}

There is an implementable finite-dimensional form when the state sets are
finite and each infimum in the selected lower operator has an exact
representation by a nonempty finite complete menu with finite real penalties.
Any piece with penalty $+\infty$ is omitted: bounded continuations and
normalization make the infimum finite, and such a piece cannot affect it.
Write the resulting entire menu at $(t,x)$ as
$q_{txj}$, $j\in\mathcal J_{tx}$, and its verified nonnegative finite
penalties as $\underline k_{txj}$. The word complete means that minimizing
over this menu gives the specified lower operator for every continuation used below;
it is not a list of favorable witnesses sampled from a larger risk set.
Let $\nu_0$ be an initial-state probability vector and put
$B_0=1$, $B_{t+1}=B_t\beta_t$.

\begin{proposition}[Lower-gain occupations]
\label{prop:gainoccupations122}
Under this finite complete-menu assumption,
$\nu_0^{\mathsf T}\underline H_0$ equals the minimum of
\begin{equation}\label{eq:gainflowobjective122}
 \sum_{t=0}^{T-1} B_t\sum_x m_t(x)b_t(x)
 +\sum_{t=0}^{T-1} B_{t+1}\sum_{x,j}
                         y_t(x,j)\underline k_{txj}
\end{equation}
over nonnegative flows satisfying
\begin{align}
 m_0&=\nu_0,\qquad
 \sum_jy_t(x,j)=m_t(x),\nonumber\\
 m_{t+1}(x')&=\sum_{x,j}y_t(x,j)q_{txj}(x').
                    \label{eq:gainflowconstraints122}
\end{align}
Every verified subsolution in Theorem~\ref{thm:dynamicgain122} gives a
lower bound $\nu_0^{\mathsf T}h_0$ on this minimum. A feasible flow alone
gives an \emph{upper} bound on the minimum and therefore does not certify
the original policy gain.
\end{proposition}

\begin{proof}
The lower recursion is a finite-horizon minimizing control problem with
risk index $j$, stage cost $b_t(x)+\beta_t\underline k_{txj}$, transition
$q_{txj}$ and zero terminal cost. Finite minima attain a minimizing Markov
choice at each date and state. Backward induction therefore gives its
value $\underline H_t$.

Any possibly history-dependent risk-index rule induces the nonnegative
occupation masses in \eqref{eq:gainflowconstraints122}. Conversely, given
a feasible flow, use the conditional probabilities
$y_t(x,j)/m_t(x)$ when $m_t(x)>0$, and any available index when the mass
is zero. Forward induction recovers exactly the given masses. Its
expected discounted cost is \eqref{eq:gainflowobjective122}, proving
the equality of minima.

Alternatively, multiply
$h_t(x)\le b_t(x)+\beta_t\{\sum_{x'}q_{txj}(x')h_{t+1}(x')+
 \underline k_{txj}\}$
by $B_ty_t(x,j)$ and sum over all $t,x,j$. The flow equalities cancel
the intermediate terms; the zero terminal value leaves
$\nu_0^{\mathsf T}h_0$ no greater than the flow objective. This proof
does not divide by a discount and also covers zero discounts. Minimizing
preserves the lower bound, whereas inserting one feasible flow can only
bound the minimum from above.
\end{proof}

These are occupations under original risk-envelope kernels, not the
physical law of the investment state. For a continuous risk set, replacing
the infimum by a subset increases it and can destroy the lower certificate.
A numerical implementation must instead bound the complete infimum from
below, possibly by a verified relaxation or an explicitly bounded
integration error. The positive penalty and the minimization direction
are both essential. Discarding nonnegative penalties is conservative
here; freezing a single reference-maximizing risk kernel need not be.

\subsection{The paid payload and the original incumbent}

To apply the action safeguard at several dates, its scalar objective must
represent the original one-step cost under the incumbent future:
\[
 Q_t^\pi(x,a)=c_t(x,a)+\beta_t\rho_{t,x,a}(V_{t+1}^\pi).
\]
Suppose, at a given state and date, one has proved a decomposition
$Q_t^\pi(x,a)=q_t(x,a)+\eta_t(x,a)$, with the strong convexity,
smoothness and paid-support assumptions of
Theorem~\ref{thm:payloadsafeguard122} for $q_t$, and an all-action bound
$0\le\eta_t(x,a)\le\delta_t(x)$. Then the realized feasible candidate
with checked displacement $d$ has the advantage lower bound
\[
 b_t(x)= -\overline m_{t,x}(d)-\delta_t(x)>0
\]
when its gain test passes. Otherwise retention gives $b_t(x)=0$.
A finite candidate menu and fixed tie conventions give a measurable
policy when the evaluated inputs are measurable. The all-state hypotheses
and any additional policy-qualification tests must also hold before
Theorem~\ref{thm:dynamicgain122} is invoked.

In the executed two-date economy, the unchanged last-date controller
provides its already proved allowance. That particular numerical allowance
cannot be reused at every date of a longer learned recursion without a
new proof. If $q_t$ uses the original optimal future while the actor uses
$V_{t+1}^\pi$, a valid nonnegative allowance follows, for example, from
$0\le V_{t+1}^\pi-V_{t+1}^*\le\epsilon_{t+1}$ uniformly and the original
monotone cash-invariant risk evaluation:
$0\le\eta_t\le\beta_t\epsilon_{t+1}$.
The continuation level and its error must therefore be retained and
propagated even when the action choice uses marginal information.
No cost or sample claim for such a longer implementation is made here.

\subsection{A stationary extension with a complete domain}

Let $\mathcal X$ be a measurable state space and let
$B_b(\mathcal X)$ be the Banach space of all bounded measurable real
functions with the supremum norm. Fix feasible measurable stationary
policies $\pi,\pi'$ and a discount $0\le\beta<1$. Assume their selected
stage costs are bounded, and their original normalized, monotone,
cash-invariant risk maps send \emph{every} member of $B_b(\mathcal X)$
to a member of $B_b(\mathcal X)$. Assume the stated dual representations
on this entire space. These complete-domain conditions are needed for
the contraction argument; a collection containing only the eventual
solutions is not a replacement for them.

Define $T_\zeta v(x)=c(x,\zeta(x))+\beta\rho_{x,\zeta(x)}(v)$ for
$\zeta\in\{\pi,\pi'\}$.

\begin{proposition}[Stationary certified policy improvement]
\label{prop:stationarygain122}
The maps $T_\pi,T_{\pi'}$ have unique fixed points
$V^\pi,V^{\pi'}\in B_b(\mathcal X)$. Form the advantage $b^*$,
centered operator $\mathcal L^\pi$, and a bounded measurable certificate
$0\le b\le b^*$ by their stationary versions above. Let
$\underline{\mathcal L}$ be a verified lower-penalty operator that maps
all of $B_b(\mathcal X)$ to $B_b(\mathcal X)$.
Then the maps
\[
 h\longmapsto b+\beta\mathcal L^\pi h,
 \qquad h\longmapsto b+\beta\underline{\mathcal L}h
\]
have unique fixed points $H,\underline H$, and
\[
       0\le\underline H\le H\le V^\pi-V^{\pi'}.
\]
Iteration of either lower-gain map from zero increases monotonically
and converges uniformly to its fixed point. Every bounded nonnegative
subsolution of the conservative map is a valid original-gain lower
certificate.
\end{proposition}

\begin{proof}
The monetary risk assumptions imply nonexpansiveness on the supremum
norm by the order and cash-invariance argument in
Lemma~\ref{lem:lowercenter122}. Thus $T_\zeta$ is a $\beta$-contraction
of the stated Banach space into itself. The contraction theorem gives
its unique fixed point. Lemma~\ref{lem:lowercenter122} makes each
gain map a contraction on the same complete space, proving existence,
uniqueness and uniform convergence of its iterates.

Normalization and $b\ge0$ give a nonnegative first iterate from zero.
Monotonicity makes all subsequent iterates nondecreasing. The exact
identity
\[
 D:=V^\pi-V^{\pi'}=b^*+\beta\mathcal L^\pi D
\]
holds by subtracting the two original fixed-point equations.
The map on the right is also a contraction. Its iterates from zero
are nonnegative because $b^*\ge b\ge0$, and converge to $D$.
Comparing the three iterations step by step, using the operator order
and $b\le b^*$, yields $0\le\underline H\le H\le D$ in the limit.
If $h\ge0$ is a conservative subsolution, iteration from $h$ is
nondecreasing and, by contraction, converges uniformly to
$\underline H$. Hence $h\le\underline H\le D$.
\end{proof}

The stationary statement proves improvement under its all-state original-
objective premises. It does not prove convergence of an arbitrary learned
policy-iteration sequence to the optimum, a complexity bound for computing
the risk infima, a neural representation advantage, or a complete-cost
adoption result. Its economic use is precise: a verified action advantage
can be propagated through the same recursive preference while preserving
the continuation information and adverse-kernel effects that a pointwise
fit or a physical expectation would miss.
